\documentclass[11pt]{article}

\usepackage[utf8]{inputenc}
\usepackage[T1]{fontenc}
\usepackage{lmodern}
\usepackage{geometry}
\usepackage{amsmath,amssymb,amsfonts,amsthm,mathtools}
\usepackage{bm}
\usepackage{enumitem}
\usepackage{microtype}
\usepackage{hyperref}
\usepackage{titlesec}
\usepackage{booktabs}
\usepackage{array}
\usepackage{setspace}
\usepackage{mathrsfs}
\usepackage{physics}

\geometry{margin=1in}
\hypersetup{colorlinks=true, linkcolor=black, urlcolor=blue, citecolor=black}
\setlist[enumerate]{topsep=0.4em,itemsep=0.35em}
\setlist[itemize]{topsep=0.3em,itemsep=0.25em}
\titleformat{\section}{\large\bfseries}{\thesection.}{0.5em}{}
\titleformat{\subsection}{\normalsize\bfseries}{\thesubsection.}{0.5em}{}
\setstretch{1.06}

\newcommand{\Aether}{\AE{}ther}
\newcommand{\Aflow}{\AE{}ther-flow}
\newcommand{\TheoryName}{\Aflow{} Interpretation of Relativity}
\newcommand{\TheoryProgram}{\Aether{} / \Aflow{} framework}

\newcommand{\Car}{\mathrm{car}}
\newcommand{\Cur}{\mathrm{cur}}
\newcommand{\Hom}{\mathrm{hom}}

\newcommand{\ForPrimSpace}[1]{\mathcal I_{#1}^{\mathrm{fpr}}}
\newcommand{\ForSpace}[1]{\mathcal I_{#1}^{\mathrm{for}}}
\newcommand{\ForPrimInc}[1]{\iota_{#1}^{\mathrm{fpr}}}

\newcommand{\Reduction}[1]{\mathcal R_{#1}}
\newcommand{\ReductionTriple}{\mathcal R_{\mathrm{tot}}}
\newcommand{\ObsShell}{\mathcal S_{\mathrm{obs}}}
\newcommand{\PrimShell}{\mathcal S_{\mathrm{prim}}}

\newcommand{\ObsCurrent}[1]{\mathcal J_{#1}}
\newcommand{\ObsTheta}[1]{\Theta_{#1}}
\newcommand{\PrimCurrent}[1]{\mathcal J_{#1}^{\mathrm{prim}}}
\newcommand{\PrimTheta}[1]{\Theta_{#1}^{\mathrm{prim}}}
\newcommand{\EqProj}[1]{P_{#1}^{\mathrm{eq}}}

\newcommand{\DefectSpace}{\mathcal D}
\newcommand{\DefectTuple}{\mathbf d}
\newcommand{\PrimDefectTuple}{\mathbf d^{\mathrm{prim}}}

\newcommand{\ObsResidualSector}{\mathfrak R_{\mathrm{res}}}
\newcommand{\ObsCompletionSector}{\mathfrak R_{\mathrm{cmp}}}
\newcommand{\ObsMatterSector}{\mathfrak R_{\mathrm{mat}}}
\newcommand{\PrimResidualSector}{\mathfrak R_{\mathrm{res}}^{\mathrm{prim}}}
\newcommand{\PrimCompletionSector}{\mathfrak R_{\mathrm{cmp}}^{\mathrm{prim}}}
\newcommand{\PrimMatterSector}{\mathfrak R_{\mathrm{mat}}^{\mathrm{prim}}}

\newcommand{\ResidualRead}{\mathscr R_{\mathrm{res}}}
\newcommand{\CompletionRead}{\mathscr R_{\mathrm{cmp}}}
\newcommand{\MatterRead}{\mathscr R_{\mathrm{mat}}}
\newcommand{\MetricOut}{\mathscr G_{\mathrm{out}}}
\newcommand{\MatterOut}{\mathscr T_{\mathrm{out}}}

\newcommand{\ResidualTarget}{\mathcal V_{\mathrm{res}}}
\newcommand{\CompletionTarget}{\mathcal V_{\mathrm{cmp}}}
\newcommand{\MatterTarget}{\mathcal V_{\mathrm{mat}}}

\newtheorem{definition}{Definition}
\newtheorem{lemma}{Lemma}
\newtheorem{proposition}{Proposition}
\newtheorem{remark}{Remark}
\newtheorem{corollary}{Corollary}
\newtheorem{theorem}{Theorem}

\title{\textbf{The \TheoryName{}}\\[0.4em]
\large Primitive-Reservoir Observer-Sector\\
\large Defect-Fiber Factorization Theorem}
\author{}
\date{}

\begin{document}

\maketitle

\begin{abstract}
After the primitive-reservoir observer-reduction datum origin theorem, the current / projected-equilibrium part of the primitive observer-reduction datum is no longer primitive at current scope. The next honest burden is therefore to prove that the observer-sector outputs themselves factor through the primitive defect fibers.

The present note proves exactly that and no more. It assumes the already fixed shell-reduction maps, the already derived primitive current parents and primitive equilibrium readouts / projectors, and the already recorded observer-side continuity-Hessian sector readouts. From those inputs it shows that the observer-sector outputs pulled back to the primitive shell are constant on fibers of the primitive six-defect map and hence admit unique factor maps on the image of that map. Equivalently, the recorded observer readouts restrict to the primitive defect image and already produce the primitive observer-sector outputs there.

The gain is narrow but front-line. This note does not yet derive the one operative metric output map \(\MetricOut\) or the ordinary matter output map \(\MatterOut\). It only removes the remaining conditionality at the defect-fiber level so that the next burden can stay on that same reduction chain rather than widening early to another deeper-origin relay.
\end{abstract}

\section{Introduction}

The active derivational program of \TheoryName{} has now reached the point where staying honest requires staying local. The primitive-reservoir observer-reduction datum origin theorem has already dealt with the current / equilibrium constituent of the datum. That gain matters because it removes the first still-primitive block inside the primitive observer-reduction chain itself. Once that is done, the next front-line question is not another same-output deeper-origin relay. It is whether the observer-sector outputs on that same chain are already forced to descend through the primitive defect fibers.

That is the sole target of the present note. It does not yet try to derive the one operative metric output map \(\MetricOut\) or the ordinary matter output map \(\MatterOut\). It instead isolates the sector-output stage that sits immediately between the already derived primitive six-defect tuple and those later metric / matter output maps. The theorem proved here is that, once the active shell reduction maps and the already recorded observer-side continuity-Hessian readouts are granted, the primitive observer-sector outputs are no longer extra datum. They are uniquely determined by the primitive defect tuple itself.

\paragraph{Framework and claim boundary.}
\TheoryProgram{} denotes the broader \Aether{} / \Aflow{} framework built on the claims that the \Aether{} is the underlying four-dimensional substrate of reality, the \Aflow{} is its intrinsic ordered motion, observed three-dimensional space is the local experiential slice of that deeper substrate, S-time is the experienced order of change arising from matter, light, and the \Aflow{}, and observed expansion is the three-dimensional appearance of deeper four-dimensional ordered motion. Within that framework, \TheoryName{} denotes the exact-closure theory in which the effective gravitational dynamics are adopted to be exactly Einsteinian with universal matter coupling. In the active sequence, \emph{adoption} means use of that established relativistic dynamics without claiming substrate derivation, while \emph{derivation} is reserved for a first-principles recovery from explicit substrate variables.

The present note stays strictly inside that boundary. It does not claim a completed first-principles derivation of GR. Its narrower claim is only that, on the active primitive-reservoir line, the observer-sector outputs of the continuity-Hessian package factor through the primitive defect fibers once the current / equilibrium datum block has already been derived.

\section{Fixed Inputs from the Active Line}

\begin{definition}[Fixed continuation data after the datum-origin note]
\label{def:fixed}
For each sector label \(\sigma\in\{\Car,\Cur,\Hom\}\), fix the following continuation data from the active primitive-reservoir line.
\begin{enumerate}
    \item A primitive-reservoir state manifold \(\ForPrimSpace{\sigma}\), a forcing shell \(\ForSpace{\sigma}\subset\ForPrimSpace{\sigma}\), its inclusion
    \[
    \ForPrimInc{\sigma}:\ForSpace{\sigma}\hookrightarrow\ForPrimSpace{\sigma},
    \]
    and a shell reduction map
    \[
    \Reduction{\sigma}:\ForSpace{\sigma}\to X_{\sigma}
    \]
    into the visible observer-sector space.
    \item The current / equilibrium output of the primitive-reservoir observer-reduction datum origin theorem, namely maps
    \[
    \PrimCurrent{\sigma}:\ForPrimSpace{\sigma}\to Y_{\sigma}^{\mathrm{cur}},
    \qquad
    \PrimTheta{\sigma}:\ForPrimSpace{\sigma}\to Y_{\sigma}^{\mathrm{eq}},
    \qquad
    \EqProj{\sigma}:Y_{\sigma}^{\mathrm{eq}}\to Y_{\sigma}^{\mathrm{eq}},
    \]
    together with the recorded observer-side parents
    \[
    \ObsCurrent{\sigma}:X_{\sigma}\to Y_{\sigma}^{\mathrm{cur}},
    \qquad
    \ObsTheta{\sigma}:X_{\sigma}\to Y_{\sigma}^{\mathrm{eq}},
    \]
    satisfying the shell identities
    \begin{equation}
    \PrimCurrent{\sigma}\circ\ForPrimInc{\sigma}
    =
    \ObsCurrent{\sigma}\circ\Reduction{\sigma},
    \qquad
    \PrimTheta{\sigma}\circ\ForPrimInc{\sigma}
    =
    \ObsTheta{\sigma}\circ\Reduction{\sigma}.
    \label{eq:shell_restriction}
    \end{equation}
    \item The observer-shell product
    \[
    \ObsShell
    \equiv
    X_{\Car}\times X_{\Cur}\times X_{\Hom},
    \]
    the defect target
    \[
    \DefectSpace
    \equiv
    Y_{\Car}^{\mathrm{cur}}
    \times
    Y_{\Cur}^{\mathrm{cur}}
    \times
    Y_{\Hom}^{\mathrm{cur}}
    \times
    Y_{\Car}^{\mathrm{eq}}
    \times
    Y_{\Cur}^{\mathrm{eq}}
    \times
    Y_{\Hom}^{\mathrm{eq}},
    \]
    and the already recorded observer defect tuple
    \begin{equation}
    \begin{aligned}
    \DefectTuple(x_{\Car},x_{\Cur},x_{\Hom})
    \equiv\;
    \Bigl(
    &\ObsCurrent{\Car}(x_{\Car}),
    \ObsCurrent{\Cur}(x_{\Cur}),
    \ObsCurrent{\Hom}(x_{\Hom}),
    \\
    &(\mathbf 1-\EqProj{\Car})\ObsTheta{\Car}(x_{\Car}),
    (\mathbf 1-\EqProj{\Cur})\ObsTheta{\Cur}(x_{\Cur}),
    \\
    &
    (\mathbf 1-\EqProj{\Hom})\ObsTheta{\Hom}(x_{\Hom})
    \Bigr)
    \end{aligned}
    \label{eq:obs_tuple}
    \end{equation}
    on \(\ObsShell\).
    \item The already recorded observer-side sector outputs
    \[
    \ObsResidualSector:\ObsShell\to\ResidualTarget,
    \qquad
    \ObsCompletionSector:\ObsShell\to\CompletionTarget,
    \qquad
    \ObsMatterSector:\ObsShell\to\MatterTarget,
    \]
    together with observer-side readout maps
    \[
    \ResidualRead:\operatorname{im}\DefectTuple\to\ResidualTarget,
    \qquad
    \CompletionRead:\operatorname{im}\DefectTuple\to\CompletionTarget,
    \qquad
    \MatterRead:\operatorname{im}\DefectTuple\to\MatterTarget,
    \]
    such that the continuity-Hessian bridge theorem already gives the observer-side factorization identities
    \begin{equation}
    \ObsResidualSector=\ResidualRead\circ\DefectTuple,\qquad
    \ObsCompletionSector=\CompletionRead\circ\DefectTuple,\qquad
    \ObsMatterSector=\MatterRead\circ\DefectTuple.
    \label{eq:obs_factorization}
    \end{equation}
\end{enumerate}
\end{definition}

\begin{remark}
The present note uses the previous datum-origin manuscript exactly as a stopping point rather than reopening it. The current parents \(\PrimCurrent{\sigma}\), equilibrium readouts \(\PrimTheta{\sigma}\), and equilibrium projectors \(\EqProj{\sigma}\) are taken here as fixed output from that note. The new burden is only to remove the next conditional block downstream of them.
\end{remark}

\begin{remark}
The matter-sector readout \(\MatterRead\) is not the same object as the later ordinary matter output map \(\MatterOut\). The former is the already recorded six-defect readout that appears in the continuity-Hessian observer-sector decomposition. The latter is the full ordinary matter output map on the primitive reduction chain and remains a later burden together with \(\MetricOut\).
\end{remark}

\section{Primitive Defect Tuple and Pulled-Back Observer Outputs}

\begin{definition}[Primitive shell, reduction product, and primitive defect tuple]
\label{def:primitive}
Define the primitive shell product by
\[
\PrimShell
\equiv
\ForSpace{\Car}\times\ForSpace{\Cur}\times\ForSpace{\Hom}.
\]
For
\[
\mathbf w=(w_{\Car},w_{\Cur},w_{\Hom})\in\PrimShell
\]
define the total shell reduction by
\[
\ReductionTriple(\mathbf w)
\equiv
\bigl(
\Reduction{\Car}(w_{\Car}),
\Reduction{\Cur}(w_{\Cur}),
\Reduction{\Hom}(w_{\Hom})
\bigr)
\in\ObsShell.
\]
Define the primitive defect tuple
\begin{equation}
\begin{aligned}
\PrimDefectTuple(\mathbf w)
\equiv\;
\Bigl(
&\PrimCurrent{\Car}(\ForPrimInc{\Car}(w_{\Car})),
\PrimCurrent{\Cur}(\ForPrimInc{\Cur}(w_{\Cur})),
\PrimCurrent{\Hom}(\ForPrimInc{\Hom}(w_{\Hom})),
\\
&(\mathbf 1-\EqProj{\Car})\PrimTheta{\Car}(\ForPrimInc{\Car}(w_{\Car})),
(\mathbf 1-\EqProj{\Cur})\PrimTheta{\Cur}(\ForPrimInc{\Cur}(w_{\Cur})),
\\
&
(\mathbf 1-\EqProj{\Hom})\PrimTheta{\Hom}(\ForPrimInc{\Hom}(w_{\Hom}))
\Bigr)
\end{aligned}
\label{eq:prim_tuple}
\end{equation}
as a map \(\PrimDefectTuple:\PrimShell\to\DefectSpace\).
\end{definition}

\begin{definition}[Primitive observer-sector outputs and primitive defect fibers]
\label{def:outputs}
Define the primitive observer-sector outputs as the pullbacks of the fixed observer-side sector outputs along the same shell reduction:
\begin{equation}
\PrimResidualSector
\equiv
\ObsResidualSector\circ\ReductionTriple,
\qquad
\PrimCompletionSector
\equiv
\ObsCompletionSector\circ\ReductionTriple,
\qquad
\PrimMatterSector
\equiv
\ObsMatterSector\circ\ReductionTriple.
\label{eq:pullback_outputs}
\end{equation}
For any
\[
\delta\in\operatorname{im}\PrimDefectTuple
\]
the corresponding \emph{primitive defect fiber} is
\[
(\PrimDefectTuple)^{-1}(\delta)
\subset\PrimShell.
\]
\end{definition}

\begin{remark}
Definition \ref{def:outputs} is the precise sense in which this note speaks of observer-sector outputs on the primitive line. They are not new observer laws. They are the already recorded observer-sector outputs viewed upstairs on the primitive shell through the same reduction chain.
\end{remark}

\section{Defect-Fiber Factorization}

\begin{proposition}[The primitive defect tuple descends to the recorded observer tuple]
\label{prop:descent}
For every \(\mathbf w\in\PrimShell\),
\begin{equation}
\PrimDefectTuple(\mathbf w)
=
\DefectTuple\bigl(\ReductionTriple(\mathbf w)\bigr).
\label{eq:descent}
\end{equation}
In particular,
\[
\operatorname{im}\PrimDefectTuple\subset\operatorname{im}\DefectTuple.
\]
\end{proposition}

\begin{proof}
Write \(\mathbf w=(w_{\Car},w_{\Cur},w_{\Hom})\). Apply the shell identities \eqref{eq:shell_restriction} sector by sector. The first three coordinates of \(\PrimDefectTuple(\mathbf w)\) become
\[
\ObsCurrent{\Car}(\Reduction{\Car}(w_{\Car})),
\qquad
\ObsCurrent{\Cur}(\Reduction{\Cur}(w_{\Cur})),
\qquad
\ObsCurrent{\Hom}(\Reduction{\Hom}(w_{\Hom})).
\]
Applying the same shell identities to the equilibrium readouts and then composing with \((\mathbf 1-\EqProj{\sigma})\) gives the remaining three coordinates
\[
(\mathbf 1-\EqProj{\Car})\ObsTheta{\Car}(\Reduction{\Car}(w_{\Car})),
\]
\[
(\mathbf 1-\EqProj{\Cur})\ObsTheta{\Cur}(\Reduction{\Cur}(w_{\Cur})),
\qquad
(\mathbf 1-\EqProj{\Hom})\ObsTheta{\Hom}(\Reduction{\Hom}(w_{\Hom})).
\]
This is exactly the right-hand side of \eqref{eq:descent} by the definition \eqref{eq:obs_tuple} of the recorded observer defect tuple.
\end{proof}

\begin{theorem}[Observer-sector outputs factor through the primitive defect fibers]
\label{thm:factor}
Assume the fixed data of Definitions \ref{def:fixed} and \ref{def:outputs}. Then the primitive observer-sector outputs satisfy
\begin{equation}
\PrimResidualSector
=
\ResidualRead\circ\PrimDefectTuple,
\qquad
\PrimCompletionSector
=
\CompletionRead\circ\PrimDefectTuple,
\qquad
\PrimMatterSector
=
\MatterRead\circ\PrimDefectTuple.
\label{eq:primitive_factorization}
\end{equation}
Consequently:
\begin{enumerate}
    \item each primitive observer-sector output is constant on every primitive defect fiber \((\PrimDefectTuple)^{-1}(\delta)\);
    \item the restrictions of \(\ResidualRead,\CompletionRead,\MatterRead\) to \(\operatorname{im}\PrimDefectTuple\) are the unique factor maps through which the three primitive observer-sector outputs descend.
\end{enumerate}
\end{theorem}

\begin{proof}
Start from the pullback identities \eqref{eq:pullback_outputs}. Using the observer-side factorizations \eqref{eq:obs_factorization} and then Proposition \ref{prop:descent},
\[
\PrimResidualSector
=
\ObsResidualSector\circ\ReductionTriple
=
\ResidualRead\circ\DefectTuple\circ\ReductionTriple
=
\ResidualRead\circ\PrimDefectTuple.
\]
The same argument gives
\[
\PrimCompletionSector
=
\CompletionRead\circ\PrimDefectTuple,
\qquad
\PrimMatterSector
=
\MatterRead\circ\PrimDefectTuple.
\]
This proves \eqref{eq:primitive_factorization}.

If \(\mathbf w,\mathbf w'\in\PrimShell\) satisfy
\[
\PrimDefectTuple(\mathbf w)=\PrimDefectTuple(\mathbf w'),
\]
then \eqref{eq:primitive_factorization} immediately yields
\[
\PrimResidualSector(\mathbf w)=\PrimResidualSector(\mathbf w'),
\]
\[
\PrimCompletionSector(\mathbf w)=\PrimCompletionSector(\mathbf w'),
\qquad
\PrimMatterSector(\mathbf w)=\PrimMatterSector(\mathbf w').
\]
Hence each output is constant on primitive defect fibers.

For uniqueness, let
\[
\widetilde{\ResidualRead}:\operatorname{im}\PrimDefectTuple\to\ResidualTarget
\]
be any map satisfying
\[
\PrimResidualSector=\widetilde{\ResidualRead}\circ\PrimDefectTuple.
\]
For every \(\delta\in\operatorname{im}\PrimDefectTuple\), choose \(\mathbf w\in\PrimShell\) with \(\PrimDefectTuple(\mathbf w)=\delta\). Then
\[
\widetilde{\ResidualRead}(\delta)
=
\PrimResidualSector(\mathbf w)
=
\ResidualRead(\delta),
\]
so \(\widetilde{\ResidualRead}\) agrees with the restriction of \(\ResidualRead\) to \(\operatorname{im}\PrimDefectTuple\). The same argument applies to the completion and matter-sector maps.
\end{proof}

\begin{corollary}[The ordered readout package is already inherited on the primitive defect image]
\label{cor:tel}
Fix
\[
\delta=(d_1,d_2,d_3,d_4,d_5,d_6)\in\operatorname{im}\PrimDefectTuple.
\]
Then the same ordered one-slot telescoping decomposition used in the continuity-Hessian bridge theorem holds on the primitive defect image. For example,
\begin{equation}
\begin{aligned}
\ResidualRead(\delta)
=\;&
\ResidualRead(d_1,0,0,0,0,0)
+[\ResidualRead(d_1,d_2,0,0,0,0)-\ResidualRead(d_1,0,0,0,0,0)]
\\
&+[\ResidualRead(d_1,d_2,d_3,0,0,0)-\ResidualRead(d_1,d_2,0,0,0,0)]
\\
&+[\ResidualRead(d_1,d_2,d_3,d_4,0,0)-\ResidualRead(d_1,d_2,d_3,0,0,0)]
\\
&+[\ResidualRead(d_1,d_2,d_3,d_4,d_5,0)-\ResidualRead(d_1,d_2,d_3,d_4,0,0)]
\\
&+[\ResidualRead(d_1,d_2,d_3,d_4,d_5,d_6)-\ResidualRead(d_1,d_2,d_3,d_4,d_5,0)].
\end{aligned}
\label{eq:tel}
\end{equation}
The same identity holds for \(\CompletionRead\) and \(\MatterRead\).
\end{corollary}

\begin{proof}
This is the standard telescoping identity in six ordered variables, applied to the restrictions of the already fixed observer-side readouts to \(\operatorname{im}\PrimDefectTuple\). The point relevant here is structural: once Theorem \ref{thm:factor} is proved, the ordered readout package already lives on the primitive defect image itself and need not be postulated there as new data.
\end{proof}

\section{Main Theorem}

\begin{theorem}[Primitive-reservoir observer-sector defect-fiber factorization theorem]
\label{thm:main}
Assume the fixed continuation data of Definition \ref{def:fixed}. Then the observer-sector output constituent of the primitive-reservoir observer-reduction datum is derived rather than separately postulated. More precisely:
\begin{enumerate}
    \item the primitive six-defect tuple \(\PrimDefectTuple\) is the pullback of the already recorded observer tuple \(\DefectTuple\) along the total shell reduction \(\ReductionTriple\);
    \item the primitive observer-sector outputs \(\PrimResidualSector,\PrimCompletionSector,\PrimMatterSector\) are the pullbacks of the already recorded observer-side sector outputs along that same reduction chain;
    \item the primitive observer-sector outputs factor uniquely through the primitive defect fibers, with factor maps given by the restrictions of the already recorded readouts \(\ResidualRead,\CompletionRead,\MatterRead\) to \(\operatorname{im}\PrimDefectTuple\);
    \item therefore the ordered primitive readout package is already determined by the primitive defect tuple and is not an extra primitive constituent of the observer-reduction datum.
\end{enumerate}
Therefore item \(2\) of the remaining front-line burden after the datum-origin theorem is now discharged at current scope. What remains open is item \(3\): derivation of the one operative metric output map \(\MetricOut\) and the ordinary matter output map \(\MatterOut\) from that same reduction chain.
\end{theorem}

\begin{proof}
Item (1) is Proposition \ref{prop:descent}. Item (2) is the definition \eqref{eq:pullback_outputs}. Item (3) is Theorem \ref{thm:factor}. Item (4) is Corollary \ref{cor:tel}. The final sentence is the routing consequence of these statements: once the primitive observer-sector outputs are forced to descend through the primitive defect fibers, the next honest burden stays on the same primitive observer-reduction chain and moves directly to the operative metric and ordinary matter output maps.
\end{proof}

\begin{corollary}[What the next primary manuscript should now target]
\label{cor:next}
After Theorem \ref{thm:main}, the next primary manuscript on the active primitive observer-reduction line should remain on the same reduction chain and should not default to another same-output deeper-origin relay. Its front-line task should be:
\begin{enumerate}
    \item derivation of the one operative metric output map \(\MetricOut\);
    \item derivation of the ordinary matter output map \(\MatterOut\) through that same one-metric chain.
\end{enumerate}
\end{corollary}

\begin{proof}
Theorem \ref{thm:main} removes the defect-fiber factorization block from the primitive observer-reduction datum. The only remaining benchmark-facing constituents named in the current routing burden are therefore the metric-output and ordinary-matter-output maps.
\end{proof}

\section{Conclusion}

The positive result of this note is narrow but real. On the active primitive-reservoir line, the observer-sector outputs are no longer treated as additional primitive data once the shell reduction maps, the already derived current / projected-equilibrium six-defect tuple, and the already recorded observer-side continuity-Hessian readouts are fixed. Pulled back to the primitive shell, those outputs are automatically constant on primitive defect fibers, so they descend uniquely to the image of the primitive defect tuple.

The scope remains disciplined. This note does not yet derive the one operative metric output map \(\MetricOut\), and it does not yet derive the ordinary matter output map \(\MatterOut\). It therefore does not complete the full primitive-reservoir observer-reduction datum. What it does do is remove the next conditional block at exactly the right level: the primitive defect-fiber stage. The next honest burden is now the metric / matter output step on that same primitive observer-reduction chain, not another deeper-origin relay that preserves the same current result.

\end{document}
